\subsection{One retained word through the whole congruence system}

\begin{statement}{E.55}{Proved}{The full preceding boundary family through every modulus prime to six}
Retain the actual source maps of E.51--E.53, with
$R=\prod_{p\geq5}\mathbf Z_p$.
Let $G_{\mathrm{old},\infty}$ be the actual closed joint
image of $\widehat B$ in the entire E.51 family: its
fixed $D_{23}$ and $G_L$ factors, its level-five and
level-seven symplectic factors, the full two-adic
class-two and three-adic class-three boundary factors,
and every full relative pro-$p$ boundary factor for
$p\geq5$. The Hurwitz source is retained through its
specified quotient map from $G_L$.
Write
$\rho_{\mathrm{old},\infty}:\widehat B\twoheadrightarrow
G_{\mathrm{old},\infty}$ for this continuous map.

Every one of the $432$ multiplier-one points, together
with every prescribed compatible full E.41 fibre and
full E.47 correction-image fibre, has a single
$f\in\widehat F'$ retaining all those images and
satisfying
\[
 \rho_R(f)=\rho_R(f_\sigma),\qquad
 \rho_L(f)=\rho_L(f_\sigma),\qquad
 \chi(\sigma)=1\in\widehat{\mathbf Z}^{\times}
 \tag{E.55.1}
\]
for the same actual Galois element $\sigma$ selected
in E.43 and retained in E.51.
Thus all eight symplectic evaluations at every
modulus prime to six, and their full inverse limit,
are the evaluations of that same Galois word.
They satisfy the pentagon, the normalized cycle
relations, the specified periodic-braid conjugacies,
and the individual braid generation requirements.
The identity shadows on all full relative pro-$p$
boundary factors remain the prescribed identities.

Here are the exact joint source and word fibres.
For $m>1$ prime to six, put $g=\gcd(m,35)$ and define
$G_1=A_2$, $\rho_1=q_{A_2}$ and $C_1=1$.
For the other values of $g$, use the source $G_g$
of E.52. All maps below are the natural reductions
of those same source words. Put
\[
 \begin{split}
 \mathcal Q_m&=A_{2m}\mathbin{\times_{A_{2g}}}G_g,\\
 K_m&=C_m\cap\ker(G_m\longrightarrow G_g).
 \end{split}
 \tag{E.55.2}
\]
For every finite old source $G_{\mathrm{old}}$ retaining
$P_{4m}=B/(\gamma_3(H)H^{4m})$, the maximum common
quotient through the specified maps is precisely
$\mathcal Q_m$, and the actual joint source is
\[
       G_{\mathrm{old}}\mathbin{\times_{\mathcal Q_m}}G_m.
 \tag{E.55.3}
\]
The corrected source supplied by $f$ is onto this
whole group. Such old sources are cofinal in the
stated E.51 family for each $m$.

Put $\mathcal A_R=\varprojlim_{(m,6)=1}A_{2m}$,
using divisibility transition maps, and put
\[
 \mathcal Q_R=\mathcal A_R\mathbin{\times_{A_{70}}}G_{35}.
 \tag{E.55.4}
\]
The full actual joint image, and the corrected image,
are both the profinite group
\[
       G_{\mathrm{old},\infty}
                   \mathbin{\times_{\mathcal Q_R}}G_R.
 \tag{E.55.5}
\]
Its new kernel over the old source is the exact
$\mathcal K_R^{35}$ of E.53.5.
All projections in these displays are continuous
epimorphisms induced by the displayed original
source maps.

For the chosen old word $f_0$ of E.51, define the
closed source subgroup
\[
 \mathscr K=R_{\mathrm{nil}}\cap\ker\rho_L
       \cap\ker\rho_5\cap\ker\rho_7
       \cap\ker\rho_{D_{23}}.
 \tag{E.55.6}
\]
The corrections $k\in\mathscr K$ in $f=f_0k$ form
the complete nonempty fibre
\[
 k_0\ker(\rho_R|_{\mathscr K}),\qquad
 \rho_R(k_0)=\rho_R(f_0)^{-1}\rho_R(f_\sigma).
 \tag{E.55.7}
\]
For the fixed full old value and fixed new Galois
value, all profinite commutator words are exactly
\[
 f\,\ker\bigl((\rho_{\mathrm{old},\infty},\rho_R)
                                          |_{\widehat F'}\bigr).
 \tag{E.55.8}
\]
These formulae retain every old parameter; a new
word is not selected separately for different
congruence moduli.
\end{statement}

\begin{proof}
We first identify each quotient map, then construct
one correction in the closed source subgroup, and
finally prove generation of the actual joint source.

\paragraph{Reduction and the whole common quotient.}
For $n\mid m$ with $m$ prime to six, reducing all
matrices in E.52's $48$ coordinates gives the actual
epimorphism $G_m\to G_n$. Its restriction
$C_m\to C_n$ is onto: in the coordinates
$C_m=\Phi_m(\operatorname{SL}_2(\mathbf Z/m),V_3(\mathbf Z/m))$,
both factors reduce surjectively.
For the matrix factor this follows by lifting the
upper and lower elementary generators, which generate
over each local factor. The additive reduction is
the ordinary coordinate reduction in the proved
free rank-three summand. The reduction to $G_1=A_2$
is the parity quotient already specified in E.52.

In particular the two level maps give an actual
isomorphism
\[
                  G_{35}\xrightarrow{\sim}G_5\times_{A_2}G_7.
 \tag{E.55.9}
\]
It is injective by CRT on each of the $48$ matrix
coordinates and on the same outer coset action.
Its image has the full kernel over $A_2$, since
E.52 identifies that kernel with
$\operatorname{SL}_2(\mathbf Z/35)\times(\mathbf Z/35)^5$.
CRT identifies this kernel with the product of
the two stated level kernels. The outer $A_2$
map is onto, proving surjectivity in
\textup{(E.55.9)} as well.
This proves that the old level-five and level-seven
maps supply the indicated $G_{35}$ map with its
original source values.

The map
\[
 G_m\longrightarrow\mathcal Q_m,
       \qquad v\longmapsto(\operatorname{ab}_{2m}(v),v\bmod g)
 \tag{E.55.10}
\]
is onto and has kernel $K_m$.
The kernel assertion is its definition, using
$\ker(G_m\to A_{2m})=C_m$ from E.52.
For surjectivity, take a pair $(a,h)$ in the
displayed fibre product and lift $a$ to $v\in G_m$.
The element $h(v\bmod g)^{-1}$ lies in $C_g$.
Lift it through the just-proved epimorphism
$C_m\to C_g$ and multiply that lift by $v$.
Its image under \textup{(E.55.10)} is $(a,h)$.
This proves the assertion through explicit
quotient maps, including the case $g=1$.

The old sources contain the required abelian
quotient. Indeed
\[
 \gamma_3(H)\subset F',\qquad H^{4m}\subset F^{2m}
 \tag{E.55.11}
\]
give an actual epimorphism $P_{4m}\to A_{2m}$.
To see that $P_{4m}$ belongs to the specified old
boundary system, the class-two Hall coordinates
of its geometric kernel are modulo $4m$ in degree
one and modulo
$\gcd(4m,\binom{4m}{2})=2m$ in degree two.
These moduli follow directly from the power
formula and the commutator formula: powers of
the basis elements impose $4m$ on each coordinate,
and a power of a product of two basis elements
also imposes $\binom{4m}{2}$ on their commutator.
Conversely these moduli annihilate every $4m$-th
power in the collected class-two group.
The two natural maps therefore give the CRT
isomorphism
\[
                    P_{4m}\xrightarrow{\sim}P_4\times_{B/H}P_m.
 \tag{E.55.12}
\]
On $H$ its coordinate maps are precisely the CRT
maps from $(4m,2m)$ to $(4,2)$ and $(m,m)$.
The odd class-two exponent-$m$ group is the product
of its prime-power Sylow quotients. Each is a
finite characteristic quotient of the corresponding
full pro-$p$ boundary factor retained in E.42.
E.41 retains $P_4$. Thus these are actual old
quotients, and adjoining $P_{4m}$ only refines
within the specified system. Adjoining any finite
list of such quotients has the same property.

Consequently $G_{\mathrm{old}}$ has the map to
$\mathcal Q_m$ supplied by its $A_{2m}$ map and
its $G_g$ map from \textup{(E.55.9)}.
The compatibility over $A_{2g}$ holds on $s,r$,
since all maps are evaluations of the same original
source. Its image is the whole $\mathcal Q_m$,
because the source composite agrees with the
surjection \textup{(E.55.10)}.

By E.53, the exact image of the closed subgroup
\textup{(E.55.6)} is
\[
                         \rho_R(\mathscr K)=\mathcal K_R^{35}.
 \tag{E.55.13}
\]
Every remaining old boundary factor kills
$R_{\mathrm{nil}}$: its restriction to $H$ is
nilpotent at each finite level. The full relative
pro-$p$ factors have finite $p$-group quotients,
each of finite nilpotency class, and hence kill
$R_{\mathrm{nil}}$ in their inverse limits too.
The separate $D_{23}$ kernel in
\textup{(E.55.6)} retains the complete old fixed
finite source. Thus
\[
              \mathscr K\subset\ker\rho_{\mathrm{old},\infty}.
 \tag{E.55.14}
\]

Reduction of E.53.5 modulo $m$ is exactly $K_m$:
at five and seven it gives the first matrix
congruence subgroup and the corresponding additive
multiple; at every other prime it gives the
entire matrix and additive factors of $C_m$.
This is exactly the kernel of reduction from
$C_m$ to $C_g$ in E.52's coordinates.
Equations \textup{(E.55.13)--(E.55.14)} therefore give
$K_m\subset\rho_m(\ker\rho_{\mathrm{old}})$.
The common map to $\mathcal Q_m$ gives the reverse
inclusion. We have proved the exact equality
\[
                     \rho_m(\ker\rho_{\mathrm{old}})=K_m.
 \tag{E.55.15}
\]
The actual joint source is subdirect. Its kernel
over the old projection is $K_m$, and its quotient
comparison is \textup{(E.55.10)}. Lifting an old
element and then multiplying by this pure kernel
gives every compatible pair, proving
\textup{(E.55.3)} and the asserted maximum common
quotient.

For the inverse limit, the old maps
$P_{4m}\to A_{2m}$ give the continuous onto map
$G_{\mathrm{old},\infty}\to\mathcal A_R$.
Surjectivity follows either from the same source
values or from nonempty closed finite fibres.
Together with \textup{(E.55.9)} it gives the actual
map to \textup{(E.55.4)}. The kernel of
$G_R\to\mathcal Q_R$ is
$C_R\cap\ker(G_R\to G_{35})=\mathcal K_R^{35}$.
The surjectivity argument for
\textup{(E.55.10)} applies using the onto map
$C_R\to C_{35}$ and the compact lift of each
$\mathcal A_R$ value.
The same kernel equality as
\textup{(E.55.15)}, using
\textup{(E.55.13)--(E.55.14)}, now proves the
actual source identity \textup{(E.55.5)}.

\paragraph{A single correction retaining every old parameter.}
Take the complete E.51 word $f_0$ for the prescribed
old point and full boundary parameters. Retain its
same actual $\sigma$ from E.43.11.
The equality $\chi(\sigma)=1$ there is in all of
$\widehat{\mathbf Z}^{\times}$; it therefore holds
at every new congruence level simultaneously.
Both $f_0$ and $f_\sigma$ belong to $\widehat F'$,
so
\[
       d=\rho_R(f_0)^{-1}\rho_R(f_\sigma)\in C_R.
 \tag{E.55.16}
\]
E.51 supplies their equality under both $\rho_5$
and $\rho_7$. Under E.52's actual reduction maps,
the element $d$ therefore lies in
$C_R\cap\ker\rho_5\cap\ker\rho_7
=\mathcal K_R^{35}$.
Lift $d$ by the continuous surjection
\textup{(E.55.13)} and set $f=f_0k$.
The subgroup $\mathscr K$ lies in $\widehat F'$
and in the full old source kernel by
\textup{(E.55.14)}. The word $f$ consequently
retains every old value and has the required
new value \textup{(E.55.1)}.
This constructs one element of the closed full
source, with every prime and precision already
included in its target.

All $48$ induced evaluations of $f$ agree with
those of $f_\sigma$. The five-strand pentagon
and the three normalized evaluations are thus
literally the corresponding Galois evaluations.
The multiplier is exactly one, including on
each transvection and each periodic braid.
The resulting map on each new braid target is
the specified quotient map precomposed with
the actual Galois braid automorphism used in
E.38 and E.50. It is onto that target and has
the stated full-twist and periodic conjugacy
identities. The same assertion holds for their
actual joint inverse-limit image: all its
components are precomposed with that one
automorphism. It also holds after adjoining
the Hurwitz component, whose source values
factor through the retained $G_L$ map.

\paragraph{Generation of the whole common source after correction.}
For a specified finite target $G$ write
$\rho_G^f:\widehat B\to G$ for the normalized
source map with generator values
$s\mapsto\rho_G(s)\rho_G(f)$ and
$r\mapsto\rho_G(f)^{-1}\rho_G(r)$.
The order relations for these values hold in
every old factor by E.51 and in every new
factor by \textup{(E.55.1)}.
They therefore give actual continuous
homomorphisms from $\widehat B$ to their
common source targets.

On $G_L,G_5,G_7,G_R$ these are the original
maps precomposed with the same normalized
Galois automorphism $\alpha_\sigma$ of
$\widehat B$. This automorphism preserves
$\widehat F$, its characteristic subgroup
$\widehat H$, and $R_{\mathrm{nil}}$.
Transporting E.53.5 by $\alpha_\sigma$ gives
\[
 \rho_R^f\bigl(R_{\mathrm{nil}}\cap\ker\rho_L^f
                \cap\ker\rho_5^f\cap\ker\rho_7^f\bigr)
                         =\mathcal K_R^{35}.
 \tag{E.55.17}
\]
For the moment the separate $D_{23}$ condition
has not been imposed in this display.

Every finite solvable quotient of
$\mathcal K_R^{35}$ has order prime to six.
Its unrestricted $\operatorname{SL}_2(\mathbf Z_p)$
factors are topologically perfect. Its remaining
matrix factors are pro-five and pro-seven, and
its additive factors have primes at least five.
This proves the assertion also for each finite
reduction $K_m$.
The image of the subgroup in
\textup{(E.55.17)} under $\rho_{D_{23}}^f$ is a
solvable group with primes only two and three.
Its actual subdirect image with $K_m$ has
trivial common quotient: a common quotient
would be solvable and its order would have
both of these disjoint sets of prime divisors.
The projection kernels therefore give the
whole direct product. Imposing
$\ker\rho_{D_{23}}^f$ preserves the full
image $K_m$.

All remaining corrected old boundary maps
still kill $R_{\mathrm{nil}}$.
For each finite factor this follows by
applying its nilpotency class to the image
of the lower central series of $\widehat H$.
Their $H$ images are in the indicated
geometric kernels because $f\in\widehat H$
and the induced map on $B/H$ is the
identity. Thus the corrected joint image
contains the entire pure $K_m$ over
$G_{\mathrm{old}}$.

The corrected projection onto
$G_{\mathrm{old}}$ is onto by E.51.
The two corrected maps to $\mathcal Q_m$
agree. On $A_{2m}$ they are the original
map, since $f\in\widehat F'$ has trivial
image there. On $G_g$ they are the same
Galois precomposition, by the level-five
and level-seven equalities and
\textup{(E.55.9)}. These are exactly the
two components of \textup{(E.55.2)}.
The corrected image therefore lies in
\textup{(E.55.3)}, projects onto its old
factor, and contains its entire kernel
$K_m$ over that factor. Multiplying by
that kernel proves onto generation of
the whole fibre product.

This argument applies at every finite
old level and every $m$ after adjoining
the specified $P_{4m}$. Projection gives
the assertion before that refinement.
Any finite collection of levels is
dominated by one such old level and by
their least common multiple $m$.
The corrected closed image is therefore
onto every finite projection of
\textup{(E.55.5)}, and compactness makes
it the entire displayed profinite group.
Its restriction to $\widehat H$ is onto
the actual geometric kernel as well:
the retained $B/H$ map is the identity,
so a preimage of that kernel already
lies in $\widehat H$.

Finally, multiplying two preimages
under $\rho_R|_{\mathscr K}$ gives
exactly \textup{(E.55.7)}. Multiplying
two words in $\widehat F'$ with the
same full old and new values gives
exactly \textup{(E.55.8)}.
Every relation and generation assertion
above depends on those specified values,
so all words in this full fibre satisfy
them. This completes the exact source,
parameter and inverse-limit comparisons.
\end{proof}
